% TOP_PROBLEM: {"id": 3039, "title": "Rainbow AP(4) in an almost equinumerous coloring", "category_id": 2, "category": {"id": 2, "name": "combinatorics", "display_name": "Combinatorics", "description": "Counting problems, graph theory, discrete structures.", "slug": "combinatorics", "order_index": 2, "created_at": "2026-07-31T15:26:25.671Z"}, "status": "open", "problem_number": "OPG-478", "set_id": 12}
% FIRST_POSTED: 2026-10-01
% ENTRY_KIND: statement_only
% UnsolvedMath ID 3039; source catalogue code OPG-478
% !TeX program = lualatex
% Definition and sources only; this file is not an LLM solution attempt.
\documentclass[11pt,a4paper]{article}
\usepackage[margin=25mm]{geometry}
\usepackage{fontspec}
\setmainfont{lmroman10-regular.otf}[BoldFont=lmroman10-bold.otf,
 ItalicFont=lmroman10-italic.otf,BoldItalicFont=lmroman10-bolditalic.otf]
\usepackage{amsmath,amssymb,mathtools}
\usepackage{unicode-math}
\setmathfont{latinmodern-math.otf}
\AtBeginDocument{\let\setminus\smallsetminus}
\usepackage{xurl,hyperref}
\hypersetup{colorlinks=true,urlcolor=blue}
\setcounter{secnumdepth}{0}
\setlength{\parindent}{0pt}
\setlength{\parskip}{5pt}
\setlength{\emergencystretch}{3em}
\begin{document}
\section*{Rainbow AP(4) in an almost equinumerous coloring}\label{3039}
{\small UnsolvedMath ID 3039 · OPG-478 · Combinatorics · OpenGarden}
\subsection{Definitions and mathematical statement}
For a prime $p\ge5$, let $\mathbb Z_p=\mathbb Z/p\mathbb Z$ with addition
modulo $p$. A four-colouring is a map $\chi:\mathbb Z_p\to\{1,2,3,4\}$,
with colour classes $C_i=\chi^{-1}(i)$. Call it almost equinumerous if
\[
 |C_i|\in\{\lfloor p/4\rfloor,\lceil p/4\rceil\}
 \qquad(1\le i\le4).
\]
Thus all classes differ in size by at most one; the sizes necessarily sum
to $p$. This is a condition on all residues, not on intervals of representatives.

A four-term arithmetic progression in $\mathbb Z_p$ is an ordered list
$a,a+d,a+2d,a+3d$ with $a\in\mathbb Z_p$ and
$d\in\mathbb Z_p\setminus\{0\}$. Its four entries are distinct because
$p\ge5$. It is rainbow when the four values of $\chi$ on these entries are
pairwise different.

Does there exist an integer $p_0$ such that for every prime $p\ge p_0$ and
every almost equinumerous four-colouring, a rainbow four-term arithmetic
progression exists? The threshold is independent of the colouring. A
negative answer requires counterexamples for arbitrarily large primes,
not merely an exceptional small prime. Progressions may wrap around the
chosen residue representatives, since the equations are in $\mathbb Z_p$.
The source asks for the prime cyclic group: substituting a composite-order
cyclic group changes the problem. No particular common difference or
particular order of the four colours is prescribed.
\subsection{Short English statement}
Must every sufficiently large prime cyclic group, split into four almost equal colour classes, contain a four-term progression using all four colours?
\subsection{Sources}
\begin{itemize}
\item[S1] ULAM.AI and the UnsolvedMath Contributors, supplied problems.json, record 3039 (OPG-478), OpenGarden collection. Catalogue identity and source transcription; CC BY 4.0.
\url{https://huggingface.co/datasets/ulamai/UnsolvedMath}.
\item[S2] Open Problem Garden, Rainbow AP(4) in an almost equinumerous coloring. Original problem and discussion; the mathematical formulation below is expanded from the supplied source record.
\url{http://www.openproblemgarden.org/op/rainbow_ap_4_in_an_almost_equinumerous_coloring}.
\end{itemize}
\end{document}
